\documentclass[11pt]{article}
\usepackage[margin=25mm]{geometry}
\usepackage{fontspec}
\setmainfont{Latin Modern Roman}
\usepackage{amsmath,amssymb,mathtools}
\usepackage{unicode-math}
\setmathfont{Latin Modern Math}
\usepackage{xurl,hyperref}
\hypersetup{colorlinks=true,urlcolor=blue}
\setcounter{secnumdepth}{0}
\setlength{\parindent}{0pt}
\setlength{\parskip}{5pt}
\setlength{\emergencystretch}{3em}
\newcommand{\sref}[2]{\hyperlink{source:#1:#2}{[S#2]}}
\newcommand{\cataloguescope}[1]{\paragraph{Recorded target.}#1}
\newcommand{\R}{\mathbb R}
\newcommand{\C}{\mathbb C}
\newcommand{\N}{\mathbb N}
\newcommand{\Z}{\mathbb Z}
\newcommand{\Q}{\mathbb Q}
\newcommand{\F}{\mathbb F}
\newcommand{\eps}{\varepsilon}
\begin{document}
\section{Physical Realizability of Quantum Computation}\label{statement}

\subsection{Definitions and mathematical statement}

\paragraph{Shared notation.}$\mathbb N=\{1,2,\ldots\}$; the standard number fields have their usual meanings. All additional conventions are stated locally.


\paragraph{Definitions and conventions.} The standard circuit model uses finitely many qubits, preparation in a known input state, a universal set of local unitary gates, and final measurements. A uniform polynomial-size circuit family represents efficient quantum computation. Physical feasible realization requires scalable devices that preserve the necessary quantum coherence, control, and measurement accuracy under a specified noise model; it is stronger than a formal circuit construction.

\paragraph{Source scope and qualifications.} This is the workshop introduction's foundational physical question. No universal engineering-resource budget, device class, or noise law is supplied. Evidence for hardness of classical simulation is not a proof of scalable physical realizability, and fault-tolerance assumptions cannot be silently treated as established experimental facts.

\paragraph{Catalogue mathematical target.}
\paragraph{Catalogue wording retained for traceability.}
Is the standard model of quantum computation feasibly realizable?

\paragraph{Saved literature context.} The supplied catalogue's assessment dated 2026-08-21 records: Fault-tolerant quantum-computation theorems establish in-principle computation below an error threshold, but scalable physical realization of the standard model remains an engineering and physics problem. The assessment is retained as dated provenance; it is not a new comprehensive literature or open-status certification.
\subsection{Short English statement}
Determine whether scalable physical systems can feasibly implement the standard quantum-computation model.
\subsection{Research attempt}
\paragraph{Model, date and investigated scope.}
OpenAI GPT-6 Astra (configured model \texttt{gpt-6-astra}), ultra reasoning; generated 11 October 2026. Derives the circuit hybrid error budget and an explicit correlated-noise failure of a simple redundancy inference.

\paragraph{Sources and literature check.}
The complete supplied definition, including source qualifications and the dated literature assessment, was reviewed. The calculation below is an elementary restricted attempt motivated by [S2]; it does not claim novelty or a fresh literature-status determination.

\paragraph{Context and investigated scope.} The complete supplied definition, including source qualifications and the dated literature assessment, was reviewed. The calculation below is an elementary restricted attempt motivated by [S2]; it does not claim novelty or a fresh literature-status determination.

\paragraph{The uncorrected circuit error budget.} Let $\mathcal U_j$ be ideal channels and $\mathcal E_j$ their implementations, with diamond norm $\|\mathcal E_j-\mathcal U_j\|_\diamond\leq\epsilon_j$. The diamond norm is the supremum trace-norm amplification after adjoining an arbitrary finite ancilla; channels have norm one. Inserting and subtracting one ideal gate at a time gives
\[
 \|\mathcal E_L\cdots\mathcal E_1-\mathcal U_L\cdots\mathcal U_1\|_\diamond
 \leq\sum_{j=1}^L\epsilon_j.
\]
Thus a direct implementation with no correction needs a per-gate bound of order $1/L$ to obtain a fixed worst-case error from this estimate. It is a sufficient error budget, not a necessary lower bound for every circuit.

\paragraph{Why a noise model is substantive.} In the elementary three-bit repetition code with majority decoding, independent bit flips of probability $p$ give logical failure $3p^2(1-p)+p^3=3p^2-2p^3$. If instead all three bits flip together with probability $p$, each physical marginal still has error $p$, but the logical failure is $p$. Hence the same marginal error rate does not certify even this simple redundancy scheme. The example is classical and is not a proof of quantum fault tolerance.

\paragraph{Full-source gap.} Formal circuits and conditional correction estimates do not establish that a scalable device satisfies the required correlation, locality, control and resource assumptions. The first missing step is an experimentally justified family of microscopic noise/control models with scalable bounds. No universal engineering budget was specified in the source, and no present-day hardware or threshold status is inferred here. The mathematical checks explain why feasibility cannot be concluded from a gate set and a single error number alone.

\paragraph{Verification and adversarial audit.}
Checked channel telescoping and norm convention, independent versus correlated majority failure, and classical versus quantum scope. Model self-audit only; no independent human proof checking. Prompts A and B were applied to this canonical record and its qualified source scope. Checked channel telescoping and norm convention, independent versus correlated majority failure, and classical versus quantum scope. Retained unresolved status and zero original-target progress; A physically justified scalable noise/control/resource model is not supplied.

Coordinator checked the channel-composition telescoping bound and the different logical error probabilities for independent and common bit flips. This additional check was model-assisted, not independent human verification.

\paragraph{Tools, process and contributions.}
Python 3 was used for file preparation and structural checks; no numerical experiment or theorem prover was used. Contribution prompt A was followed by prompt B. GPT-6 Astra with ultra reasoning drafted the research and performed its stated self-audit; any additional coordinator contribution is identified in the verification and audit. A GPT-6-based Codex coordinator assembled the manuscript and checked repository integration; its exact variant and reasoning configuration were not exposed. No human mathematical contribution, independent human proof check or formal proof certification is supplied. Complete derivations and the consulted source references are included here; private reasoning transcripts are not published.

\paragraph{Final status and first remaining gap.}
\textbf{UNRESOLVED}. A physically justified scalable noise/control/resource model is not supplied.

\section{Completion Estimate}
\noindent\textbf{COMPLETION: 0\%}
The calculation is an auxiliary or restricted attempt and establishes no newly verified portion of the full source target. A physically justified scalable noise/control/resource model is not supplied.

\subsection{Sources}
\begin{itemize}
\item[\hypertarget{source:statement:1}{S1}] ULAM.AI and the UnsolvedMath Contributors, UnsolvedMath, catalogue entry. Catalogue curation under CC BY 4.0; definitions and English exposition added locally.
\url{https://huggingface.co/datasets/ulamai/UnsolvedMath}.
\item[\hypertarget{source:statement:2}{S2}] Complexity Theory (2009), Oberwolfach Reports, DOI 10.4171/owr/2009/52, pp. 2787--2850; official report source and mathematical context.
\url{https://doi.org/10.4171/owr/2009/52}.
\item[\hypertarget{source:statement:3}{S3}] Official EMS Press PDF of the same Oberwolfach report; PDF page 2.
\url{https://ems.press/content/serial-article-files/46253}.
\item[\hypertarget{source:statement:4}{S4}] fault-tolerant quantum computation theorem; reference in the saved dated catalogue literature review.
\url{https://arxiv.org/abs/quant-ph/9611025}.
\end{itemize}
\label{end:statement}
\end{document}
